\section{Evaluation protocol}
\label{sec:protocol}

\subsection{Outcomes, user model and baselines}
The primary safety outcome is the wrong-commit rate, always reported with the commit rate and the error among commits. Speed is secondary. We simulate a user who reads the output, accepts or corrects a commit, picks from an offer, or types the sentence with a fallback keyboard. For item $i$ with intended sentence $y_i$, the simulation returns the total time $t_i$ and the final text, and
\begin{equation}
  \text{\ewpm{}} = 60 \cdot \frac{\sum_i \mathbf{1}[\text{final text}_i = y_i]\;\mathrm{words}(y_i)}{\sum_i t_i}.
\end{equation}
Like the BCI utility metric \citep{dalseno2010utility}, \ewpm{} charges the time needed to fix errors. A \emph{silent error} is a wrong final text that the user did not correct; we report its rate separately.

The main user types with a gaze keyboard at 20 characters per minute and scans and selects one of $k = 3$ word alternatives in 4.5\,s. The typing speed was measured in one eye-gaze user with ALS, and the 4.5\,s are close to the 4.7\,s measured in the same study \citep{cai2024gaze}; all other costs are assumptions (\cref{tab:user}). After a \act{commit}, the user misses each wrong word with probability 0.05; after an \act{offer}, the main user sees every error. This asymmetry favours offers and is not measured. The neural input runs at 159.3 (intracortical), 75.8 (sEMG) and 152 (MEG) characters per minute. The intracortical rate is the number of characters of the cue sentences divided by the sum of the \code{speech\_duration\_s} field of the data file. The file does not say whether this field includes the 6\,s finalisation pause of the original interface, and we add no time for that pause, which favours the neural input. Because only two costs are measured, we refit every policy under eight user-model variants (\cref{app:user}).

\emph{Baselines:} always commit the top candidate; always offer and never commit; the \emph{same fitted policy without commits} (the same fit-A pipeline and policy with $\lambda = +\infty$, so every item takes the policy's non-commit path), which is the paired comparator for the speed of a commit rule; the plain threshold; gaze typing alone (error-free, which favours it); and gaze typing with \emph{word completion}, where the keyboard shows the three most frequent words with the typed prefix (a context-free unigram).

\subsection{Datasets}
\paragraph{Intracortical speech (Brain-to-Text '25).} The public copy-task file holds the online phoneme CTC logits, cue sentences and published LM-rescored sentences of participant T15 of the BrainGate2 clinical trial \citep{card2024accurate,card2025data}. We use the 1,628 trials with the 125,000-word vocabulary (14 post-implant days, 30--208 trials each). The file has no time stamps, so the file order within each day is the trial order. We evaluate our phoneme-to-word decoder (test WER 0.149) and the published sentence (WER 0.043), gated as in \cref{sec:policy}. A trial is ``wrong'' if the top sentence differs from the cue; the file has no per-trial participant ratings, so we cannot check this label against the participant's own judgement.

\paragraph{Forearm sEMG and MEG typing.} The Toronto sEMG data hold isolated key presses from 19 participants over two days \citep{achituve2025emg,achituve2024data}. We compose held-out Tatoeba sentences from single presses with oracle timing and no inter-key context (570 fit, 551 test sentences), under three decoder conditions: \emph{session-calibrated}, \emph{cross-session} and \emph{cross-subject}. In the SpanishBCBL MEG data, healthy participants typed memorised Spanish sentences \citep{levy2026brain2qwerty,spanishbcbl}; our asynchronous decoder gets no key times at inference, and the test set has 637 instances of 17 sentences (\cref{app:meg}). We use sEMG and MEG to test the stream format and the safety of the layer on weak decoders, not to test the commit rule across decoders.

\subsection{Calibration data and deployment tests}
\label{sec:splits}
The \emph{deployment protocol} calibrates on data that played no other role: a fit-A part fits the temperatures, conformal sets, guard and policy, and a disjoint calibration part only calibrates.

\paragraph{Fixed splits.} On Brain-to-Text '25 the pipeline is fitted on the first half of each day (812 trials) and tested on the second half (816). Each day's fit trials are split into fit-A (411 trials) and clean calibration (401), by alternating trials (split A) or by taking the earlier half as fit-A (split B). On sEMG, 142 fit sentences fit the policy and 143 calibrate; on MEG, 6 validation sentences fit and 5 calibrate.

\paragraph{Time-ordered (forward-chaining) tests.} A real deployment calibrates on earlier data and runs on later data. On Brain-to-Text '25 we therefore fit, calibrate and test in day order: (V1) fit on days $1..k{-}2$, calibrate on day $k{-}1$, test on day $k$; (V4) fit on days $1..k{-}3$, calibrate on days $k{-}2$ and $k{-}1$; (V2) fit once on days 1--3, calibrate on all earlier non-fit days; (V3, session start) fit on days $1..k{-}1$, calibrate on the first half of day $k$ and test on its second half. These give 10--13 test days per design. sEMG has no session order (sentences are composed from presses drawn at random), so it has no time-ordered test.

\paragraph{Random re-splits.} We also repeat the protocol 50 times with random assignments. On Brain-to-Text '25, each re-split permutes every day's trials and splits them day-stratified into fit-A (1/4), clean calibration (1/4) and test (1/2); on sEMG, each participant's fit sentences are re-split into fit-A, policy and calibration parts, with the fixed 551 test sentences, so the sEMG spread is fit and calibration variation only. To cache decodes, the fitted frame temperature was rounded to a step of 0.02 (Brain-to-Text '25) or 0.01 (sEMG); the largest relative change was 0.27\% (Brain-to-Text '25) and 0.38\% (sEMG). Calibration and test items are exchangeable by construction, so these re-splits are a Monte-Carlo check of the theorems on one pool of data, not a deployment test. Their spread is split-to-split variation, not a CI for new data. We also report the \emph{population} exceedance: we draw calibration sets with replacement from a split's pool, so that the pool is the exact population, and evaluate the chosen threshold on the whole pool (\cref{app:population}).

\subsection{Controls and statistics}
\emph{Signal-destroying} controls keep the stream format but remove the neural information (frame-shuffled and near-uniform streams, zero-input output, a permuted-label decoder); a stream-level layer should not commit these. The \emph{intent-mismatch} control is the real stream of a different sentence by the same participant, which a stream-level layer cannot detect.

All bootstrap CIs are 95\% percentile intervals from one bootstrap ($B = 2000$, seed 20261112) at the primary unit: days on Brain-to-Text '25 (14 clusters), participants on sEMG (19), sentences and participants at once on MEG. They are conditional on the fitted pipeline and threshold. Because 14 clusters are few, we add exact day-level sign tests for key paired differences. Single rates and shares of splits or days get exact Clopper--Pearson (CP) 95\% intervals; rates with no events get one-sided 95\% CP upper bounds (items treated as independent, which is optimistic under clustering). The primary endpoints were named in revision 1, after the first evaluation runs; the deployment protocol, the re-splits, the time-ordered tests, the selective controllers and the choice of the primary rule came in later revisions. All inference is therefore exploratory. All decoder settings, user-model values, the policy grid, $\alpha$, $\beta$ and $\delta$ were fixed before the runs that use them.
